\subsection{Dados e representação em grafos}
\label{sec:dados_grafos}

Os vértices dos nomes de arquivo não estavam no referencial dos sensores, e os \num{1152} canais das tampas tinham $y$ e $z$ com o sinal trocado; as duas correções foram aplicadas na construção da base de dados e conferidas nos \num{9364} eventos. Três canais quentes, com contagem em todos os eventos e \qty{47.1}{\percent} das contagens brutas, foram substituídos pela média dos vizinhos sãos: a base guarda a luz já corrigida. Uma posição foi simulada duas vezes, e três arquivos eram cópias de outros. Depois das correções, o centroide da luz detectada acompanha o vértice, com coeficiente de correlação de Pearson $r$ de \num{0.908} em $x$, \num{0.871} em $y$ e \num{0.999} em $z$ nos \num{9364} eventos.

Os eventos foram divididos em \num{7490} de treino, \num{937} de validação e \num{937} de teste, estratificados em \num{9} faixas por eixo (\num{729} grupos), divisão comum a todos os modelos. Pela simetria dos canais sob trocas de sinal dos eixos, cada evento de treino foi espelhado nos outros octantes: \num{56884} linhas de treino, \num{7.59} por evento; os de validação e de teste ficam no octante simulado. Os vértices de teste distam, em média, \qty{49.31}{\cm} do vértice de treino mais próximo e, salvo a posição simulada duas vezes, são posições da grade ausentes do treino: os erros medem interpolação numa grade de \qty{0.5}{\m}, com um sorteio de fótons por evento.

Os nós do grafo $G = (\mathcal{V}, \mathcal{E})$ são os $|A|$ canais ativos do evento, \num{245} em média nos eventos de teste, e um nó global \cite{gilmer2017neural} no centroide da luz. Cada canal recebe arestas dos $k = {}$\num{4} canais ativos mais próximos, num grafo de $k$ vizinhos mais próximos (\mbox{$k$-NN}), e troca uma aresta em cada sentido com o nó global: seis arestas dirigidas por canal (\num{248} nós e \num{1482} arestas no evento da Figura~\ref{fig:evento_exemplo}). Cada nó $i$ tem \num{17} atributos $\mathbf{f}_i$: posição, luz, fração da luz, distância e direção ao centroide, direção ao canal mais claro, luz dos vizinhos, autovalores da matriz de covariância das posições ponderada pela luz (comuns a todos os nós do evento) e posto da luz. Cada aresta tem \num{23} atributos $\mathbf{e}_{ji}$: distância, direção, inverso do quadrado da distância, diferença e semelhança de luz, e \num{16} funções de base radial da distância até \qty{7}{\m}. Os \num{33} escalares do evento, $\mathbf{u}$, são centroides com peso linear e quadrático, luz total, número de canais ativos, maior luz de canal, dispersões, assimetrias e quantis da posição. As normalizações usam só o treino (\num{8192} grafos amostrados), e cada coordenada do vértice é padronizada.
